% ===================================================================== Chapter 14
\chapter{Support Sensitivity}\label{ch:supports}

\section{Importance of End Restraint}
The end supports enter the problem twice. Axially, they decide whether the thermal growth is free (LC1) or prevented (LC2), and hence whether
a compressive force exists at all. Laterally and rotationally, they set the effective length of the resulting column and therefore its
buckling load. The real flange of the duct is not defined in this project, so three idealised scenarios with the same axial
restraint but different lateral and rotational restraint were analysed at the baseline design (Table~\ref{tab:support}). They are listed
in the order in which they were defined; they are not a ranking.

%% generated by gen_tables.py - RE-ANALYSIS 2026
\begin{table}[htbp]
\centering
\footnotesize
\caption{Support sensitivity at the baseline design (listed in order of definition, not ranked)}
\label{tab:support}
\begin{tabular}{lrrrl>{\raggedright\arraybackslash}p{0.20\linewidth}}
\toprule
Scenario & Max VM [MPa] & $\lambda_1$ & $P_{cr}$ [kN] & Mode & Mechanism first reached \\
\midrule
S1 guided ($K$ = 1) & 605.161 & 1.10805 & 608.25 & guided sway & elastic bifurcation (1.108 $<$ 1.730) \\
S2 clamped--clamped ($K$ = 0.5) & 605.160 & 4.29995 & 2,360.40 & clamped column & first yield (1.730 $<$ 4.300) \\
S3 clamped--pinned ($K$ = 0.699) & 605.160 & 2.23224 & 1,225.36 & fixed--pinned & first yield (1.730 $<$ 2.232) \\
\bottomrule
\end{tabular}
\par\smallskip\parbox{\linewidth}{\footnotesize\raggedright Source: \texttt{MASTER\_PROJECT\_DATA.csv} rows M090, M098--M110. End reaction 548,936.6~N and mean axial stress $-$582.44~MPa in all three scenarios. The mechanism column compares $\lambda_1$ with the first-yield factor of the idealised model only.}
\end{table}


\section{S1}
S1 is the LC2 support of the baseline: both end faces are held axially over their full area, which keeps them plane and prevents end rotation,
while the ends remain free to translate sideways. Structurally this is a guided (sway) column with $K = 1$. Its first mode is the sway mode of
Figure~\ref{fig:bkmode} with $\lambda_1$ = 1.108 and $P_{cr}$ = 608.25~kN.

\section{S2}
S2 adds lateral restraint at both ends ($U_\theta = 0$ on every end-face node): both ends are clamped against translation and rotation
($K = 0.5$). The first mode becomes the clamped-column shape with its maximum at mid-span (Figure~\ref{fig:s2s3}a), and $\lambda_1$ rises to
4.300 ($P_{cr}$ = 2,360.40~kN). The corresponding elastic buckling stress would be about 2500~MPa, far above yield, so a laterally held duct
would reach first yield or inelastic buckling long before this elastic bifurcation (Johnson estimate 1.58, squash 1.755).

\section{S3}
S3 clamps the inlet face and pins the outlet: the outlet face is tied to a deformable remote point on the axis that is held in position but
free to rotate ($K = 0.699$). This was the first use of a remote-point support in the project, so the exact formulation written by Mechanical
(pilot node with target element, contact elements on the outlet faces, bonded multipoint constraint) was first copied into the constant-property
toy tube (benchmark B03). The benchmark passed in two runs: the reactions matched $E\alpha\Delta T A$, no rigid-body motion occurred, and the toy
$\lambda_1$ of 2.2377 fell between the plain (2.2908) and shear-corrected (2.2256) fixed--pinned Euler values. The real S3 case then gave
$\lambda_1$ = 2.232 ($P_{cr}$ = 1,225.36~kN) with the fixed--pinned mode shape, maximum lateral deflection at 0.605 $L$ from the clamped end
(Figure~\ref{fig:s2s3}b).

\begin{figure}[htbp]
\centering
\begin{subfigure}[t]{0.49\linewidth}\centering\includegraphics[width=\linewidth,height=52mm,keepaspectratio]{LC2NS_Linear_Buckling_Mode_1_Total_Deformation_iso.png}\caption{S2, both ends clamped (load multiplier 4.2999)}\end{subfigure}\hfill
\begin{subfigure}[t]{0.49\linewidth}\centering\includegraphics[width=\linewidth,height=52mm,keepaspectratio]{S3_BK_Mode_1_iso.png}\caption{S3, clamped--pinned (load multiplier 2.2322)}\end{subfigure}
\caption{First buckling modes of the support-sensitivity scenarios (normalised eigenvectors).}
\label{fig:s2s3}
\src{\provmech}
\end{figure}

\section{Buckling-Factor Sensitivity}
The static state is identical in all three scenarios to better than 0.01\,\%: peak von Mises 605.16~MPa, mean axial stress $-$582.44~MPa and
end force 548,936.6~N (Figure~\ref{fig:supp}b). The stability, by contrast, changes by a factor of 3.9: the FE ratios S1 : S3 : S2 are
1 : 2.015 : 3.881, close to the Euler ratios 1 : 2.047 : 4 for $K$ = 1, 0.699 and 0.5 (Figure~\ref{fig:supp}a). Greater lateral and rotational
restraint of the ends shortens the effective length and raises the idealised buckling factor. The governing mechanism changes with it: under
S1 the elastic bifurcation comes before first yield (1.108 $<$ 1.730); under S2 and S3 first yield comes first (1.730 $<$ 2.232 and 4.300).

\begin{figure}[htbp]
\centering
\begin{subfigure}[t]{0.49\linewidth}\centering\includegraphics[width=\linewidth,height=52mm,keepaspectratio]{F14_support_vs_lambda1.png}\caption{First buckling load factor $\lambda_1$}\end{subfigure}\hfill
\begin{subfigure}[t]{0.49\linewidth}\centering\includegraphics[width=\linewidth,height=52mm,keepaspectratio]{F15_support_vs_LC2_stress.png}\caption{Maximum von Mises stress of the static state [MPa]}\end{subfigure}
\caption[Support sensitivity at the baseline design]{Support sensitivity at the baseline design: stability depends strongly on the end restraint, the static stress does not.}
\label{fig:supp}
\src{\provdata}
\end{figure}

\section{Implications of Support Uncertainty}
The end-support definition is the largest single uncertainty in the structural conclusion. A beam-column estimate calibrated to the FE
results (rotation-fixed ends) shows the relative lateral stiffness between the ends that would be needed: about 0.45~kN/mm for $\lambda_1$ =
1.5, about 0.71~kN/mm to reach the first-yield factor of 1.73, and about 3.9~kN/mm for full no-sway behaviour. Conversely, rotational
flexibility of real flanges would lower $\lambda_1$ below the S1 value (a conceptual $K$ = 2 gives about 0.28). None of S1, S2 or S3 is
therefore a bound on a real installation, and none is recommended here: the scenarios show the model's response to a boundary condition that
must be defined from the real mounting design before any statement about the stability of a component can be made (open task T-034).
